% Website redraw of:
%   Kristian Ranestad, Anna Seigal, Kexin Wang,
%   "A Real Generalized Trisecant Trichotomy", J. Algebra 694 (2026), 703-729,
%   Fig. 1. arXiv:2409.01356.
%
% The arXiv source includes that figure as a PNG
% (edge_quartic_dual_curve_labeled_larger_label.png), not TikZ/PGFPlots. This
% is an original PGFPlots redraw of the same mathematics: the real dual curve
% of the Edge quartic
%     25(x^4 + y^4 + z^4) - 34(x^2y^2 + x^2z^2 + y^2z^2) = 0
% in the chart w = 1 of the dual plane, with each region of the complement
% labelled by the number of real intersection points of the corresponding
% line with the quartic. The curve data (edge-dual.dat) is computed exactly
% from tangent lines by make_data.py, which also verifies every label below.
% See the publication for the original figure's rights.

\documentclass[border=4pt]{standalone}
\usepackage{pgfplots}
\pgfplotsset{compat=1.18}

\definecolor{paper}{HTML}{F7F5F0}
\definecolor{ink}{HTML}{18202A}
\definecolor{muted}{HTML}{65707A}
\definecolor{line}{HTML}{D9D6CE}
\definecolor{accent}{HTML}{1C5A76}
\definecolor{accentdark}{HTML}{123D51}
\definecolor{count4}{HTML}{B5473A}

% region labels, by number of real intersection points
\newcommand{\zero}[2]{\node[text=muted] at (axis cs:#1,#2) {0};}
\newcommand{\two}[2]{\node[text=accent] at (axis cs:#1,#2) {2};}
\newcommand{\four}[2]{\node[text=count4] at (axis cs:#1,#2) {4};}

\begin{document}
\begin{tikzpicture}[font=\sffamily\footnotesize]
\begin{axis}[
    width=6.4cm, height=6.4cm, scale only axis=false,
    xmin=-3.5, xmax=3.5, ymin=-3.5, ymax=3.5,
    axis equal image,
    axis line style={line, line width=.6pt},
    ticks=none,
    clip=true,
  ]
  \addplot[ink, line width=.85pt, unbounded coords=jump, line join=round] table {edge-dual.dat};

  % one D4-orbit per row (the figure is symmetric under the dihedral group of the square)
  \four{2.8}{2.8} \four{-2.8}{2.8} \four{2.8}{-2.8} \four{-2.8}{-2.8}
  \zero{0}{2.8} \zero{0}{-2.8} \zero{2.8}{0} \zero{-2.8}{0}
  \two{1.3}{2.5} \two{-1.3}{2.5} \two{1.3}{-2.5} \two{-1.3}{-2.5}
  \two{2.5}{1.3} \two{-2.5}{1.3} \two{2.5}{-1.3} \two{-2.5}{-1.3}
  \four{0}{1.3} \four{0}{-1.3} \four{1.3}{0} \four{-1.3}{0}
  \zero{1.0}{1.1} \zero{-1.0}{1.1} \zero{1.0}{-1.1} \zero{-1.0}{-1.1}
  \two{0.5}{0.5} \two{-0.5}{0.5} \two{0.5}{-0.5} \two{-0.5}{-0.5}
  \zero{0}{0}
\end{axis}
\end{tikzpicture}
\end{document}
